\subsection{Neighbourhoods of 0 in a metrisable topological vector space}

We say that a topological vector space E is metrisable if its topology is metrisable. Relative to the structure of its additive group and of its topology, E is, therefore, a metrisable group (GT, IX, § 3.1).

We know that, for a topological group to be metrisable, it is necessary and sufficient that there exists an enumerable fundamental system of neighbourhoods of the neutral element e, whose intersection is the single element e (GT, IX, § 3.1, prop. 1).

Also we know that the uniform structure of a metrisable topological vector space E, can be defined by an invariant distance \(d(x, y) = |x - y|\) , where \(x \mapsto |x|\) is a continuous mapping of E in \(R_{+}\) which satisfies the conditions: 1) \(|-x| = |x|\) ; 2) \(|x + y| \leqslant |x| + |y|\) ; 3) the relation \(|x| = 0\) is equivalent to x = 0 (GT, IX, § 3.1, prop. 3).

We saw (GT, IX, § 3.1, prop. 2) how such a distance \(d\) could be defined using a decreasing sequence \((\mathbf{W}_n)\) of neighbourhoods of 0 in E, forming a fundamental system of neighbourhoods and such that \(\mathbf{W}_{n+1} + \mathbf{W}_{n+1} + \mathbf{W}_{n+1} \subset \mathbf{W}_n\) . When E is a metrisable vector space over a non-discrete valued division ring K, we can also suppose that the \(\mathbf{W}_n\) are balanced (I, p.~7, prop. 4); if we revert to the process of definition of \(d\) (loc. cit.) we can see that the relation \(|\lambda| \leqslant 1\) implies that \(|\lambda x| \leqslant |x|\) . Further the conditions (EVT \(_i\) ) and (EVT \(_ii\) ) of I, p.~2 imply both that \(|\lambda x_0|\) tends to 0 as \(\lambda\) tends to 0 in K for every \(x_0 \in \mathrm{E}\) , and that \(|\lambda_0 x|\) tends to 0 as \(|x|\) tends to 0 for every \(\lambda_0 \in \mathrm{K}\) . Conversely, if the function \(|x|\) possesses all the preceding properties and if \(\mathbf{W}_n\) is the set of \(x \in \mathrm{E}\) such that \(|x| \leqslant 2^{-n}\) , then the \(\mathbf{W}_n\) form a fundamental system of balanced neighbourhoods of 0 for a metrisable topology on E that is compatible with the vector space structure of E.

Remark. --- One of the most important classes of metrisable vector spaces are the normed spaces (I, p.~3). But it must be noted that there exist metrisable vector spaces whose topology cannot be defined by a norm (I, § 3, exerc. 1); we shall study important examples later.

